\chapter{Basis-Adaptive Prime Cutoffs}
\label{ch:adaptive-cutoff}

\Cref{ch:prime-tail-certificate} proves cutoff removal for every fixed finite
Hermite section.  A different question appears when the basis size grows:
can one choose the cutoff together with the section size so that the certified
omitted tail still tends to zero?  This chapter answers that diagonal question
in the affirmative.

The result is deliberately narrower than fixed-cutoff uniformity.  The cutoff
is allowed to depend on the basis size.  No claim is made that one finite
cutoff controls every Hermite mode simultaneously.

\section{Adaptive schedule}

Let $N\ge1$ be the size of the Hermite principal section and write
\[
 D_N=2N-2
\]
for the largest polynomial degree appearing in its product kernels.  Fix a
base cutoff $Q_0\ge3$ and a slope $c>0$.  Define
\[
 U_N=\log Q_N:=\max\{\log Q_0,cN\},
 \qquad
 Q_N=\left\lceil e^{U_N}\right\rceil.
\]
When the maximum is attained by the base term, $Q_N$ is understood to be
exactly $Q_0$; this avoids introducing a floating-point successor of an
integer cutoff.

The reciprocal compactification from \Cref{ch:prime-tail-certificate} remains
available at every section:
\[
 z_{{\rm t},N}=\frac1{\log x},
 \qquad
 0\le z_{{\rm t},N}\le\frac1{U_N}.
\]
Thus increasing the basis size shrinks the compactified tail interval toward
its flat endpoint.

\section{An elementary integral envelope}

For a degree $d\ge0$, set
\[
 h_d(u)
 =
 \frac{u^{d+1}}{w^d}
 \exp\!\left(-\frac{u^2}{4w^2}+\frac u2\right),
 \qquad u>0,
\]
and define its positive logarithmic decay rate
\[
 \alpha_d(u)
 =
 \frac{u}{2w^2}-\frac12-\frac{d+1}{u}.
\]
The identity
\[
 -\frac{d}{du}\log h_d(u)=\alpha_d(u)
\]
shows that $h_d$ is decreasing once $\alpha_d$ is positive.  Moreover,
\[
 \alpha_d'(u)
 =
 \frac1{2w^2}+\frac{d+1}{u^2}>0.
\]

\begin{lemma}[Elementary adaptive tail envelope]
\label{lem:adaptive-elementary-tail}
If $\alpha_d(U)>0$, then
\[
 \int_U^\infty h_d(u)\,du
 \le
 \frac{h_d(U)}{\alpha_d(U)}.
\]
\end{lemma}

\begin{proof}
For $u\ge U$, monotonicity of $\alpha_d$ gives
$\alpha_d(u)\ge\alpha_d(U)$.  Hence
\[
 h_d'(u)=-\alpha_d(u)h_d(u)
 \le -\alpha_d(U)h_d(u).
\]
Equivalently,
\[
 h_d(u)
 \le
 h_d(U)e^{-\alpha_d(U)(u-U)}.
\]
Integration over $[U,\infty)$ gives the result.
\end{proof}

This estimate is weaker than the exact incomplete-gamma expression of
\Cref{ch:prime-tail-certificate}, but its dependence on $d$ and $U$ is
transparent enough for an asymptotic proof.

\section{A coarse Hermite-section envelope}

For $0\le m,n<N$, the Hermite linearization contains at most $N$ terms.  Each
coefficient obeys the elementary estimate
\[
 2^k k!\binom{m}{k}\binom{n}{k}
 \le
 2^{3N-3}(N-1)!.
\]
The normalized Fourier prefactor is at most one for the declared fixed width,
and the operator norm is bounded by the maximum row sum.  Therefore it is
sufficient to control the following deliberately coarse quantity:
\[
 \mathcal B_N(U)
 :=
 \frac{N^2\,2^{3N-3}(N-1)!}{\pi}
 \frac{U^{2N-1}}{w^{2N-2}}
 \frac{
 \exp\!\left(-U^2/(4w^2)+U/2\right)
 }{
 \alpha_{2N-2}(U)
 }.
\]

\begin{proposition}[Finite-section coarse majorant]
\label{prop:adaptive-coarse-envelope}
Whenever $\alpha_{2N-2}(U)>0$, the omitted prime-power tail of the
$N\times N$ Hermite section satisfies
\[
 \|E_N(e^U)\|_2\le\mathcal B_N(U).
\]
\end{proposition}

\begin{proof}
Apply \Cref{lem:adaptive-elementary-tail} to every logarithmic degree in each
entrywise certificate of \Cref{ch:prime-tail-certificate}.  Replace the number
of Hermite terms, each coefficient and every degree by the declared upper
bounds.  Finally sum at most $N$ entry bounds in each of $N$ rows.  All
replacements are one-sided and non-negative, so the resulting expression is
$\mathcal B_N(U)$.
\end{proof}

The purpose of this majorant is not numerical sharpness.  It isolates the
quadratic Gaussian term that dominates the growth of the Hermite coefficients.

\section{Exponential schedules defeat basis growth}

\begin{theorem}[Adaptive cutoff collapse]
\label{thm:adaptive-cutoff-collapse}
Fix $w>0$ and $c>0$.  With $U_N=cN$,
\[
 \log\mathcal B_N(U_N)
 =
 -\frac{c^2}{4w^2}N^2+O(N\log N).
\]
Consequently,
\[
 \mathcal B_N(cN)\longrightarrow0
 \qquad (N\to\infty).
\]
Thus the basis-adaptive schedule $Q_N=e^{cN}$ forces the certified coarse
prime-tail envelope to zero for every positive slope $c$.
\end{theorem}

\begin{proof}
Substituting $U_N=cN$ gives
\[
\begin{aligned}
 \log\mathcal B_N(cN)
 ={}&2\log N+(3N-3)\log2+\log((N-1)!)-\log\pi\\
 &+(2N-1)\log(cN)-(2N-2)\log w\\
 &-\frac{c^2}{4w^2}N^2+\frac c2N
 -\log\alpha_{2N-2}(cN).
\end{aligned}
\]
By Stirling's formula,
$\log((N-1)!)=N\log N+O(N)$.  Every displayed term except the Gaussian one
is $O(N\log N)$.  Furthermore,
\[
 \alpha_{2N-2}(cN)
 =
 \frac{cN}{2w^2}-\frac12-\frac{2N-1}{cN}
\]
is positive for all sufficiently large $N$ and contributes only $O(\log N)$
after taking a logarithm.  The negative quadratic term therefore dominates,
which proves both statements.
\end{proof}

\begin{remark}[What this theorem does not say]
The theorem controls a diagonal sequence $(N,Q_N)$.  It does not prove that a
single fixed $Q$ controls all $N$, nor that the infinite-cutoff matrices are
positive.  Those are separate problems.
\end{remark}

\section{Native PhaseNav audit}

The authoritative programme is
\[
 \texttt{construction/phasenav/secret\_of\_half\_weil\_adaptive\_cutoff\_schedule.pnv}.
\]
It declares
\[
 w=0.8,
 \qquad
 Q_0=100000,
 \qquad
 c=2,
 \qquad
 1\le N\le20.
\]
The executor evaluates the sharp v0.5 incomplete-gamma certificate at each
scheduled cutoff, while independently recording the coarse asymptotic
envelope.

For this declared schedule all twenty sharp certificates satisfy
\[
 \|E_N(Q_N)\|_2<10^{-12}.
\]
The maximum occurs before the exponential branch dominates:
\[
 \boxed{
 \max_{1\le N\le20}\|E_N(Q_N)\|_2
 \le
 3.280365246530569\times10^{-14}
 }
\]
at $N=5$.  From $N=7$ onward the bound decreases rapidly; at $N=20$ it is
approximately
\[
 6.58\times10^{-222}.
\]

The numerical evaluation is a finite audit.  The asymptotic limit itself is
provided by \Cref{thm:adaptive-cutoff-collapse}, not inferred from the twenty
samples.

\section{Claim ledger}

\begin{description}
\item[SOH-L021, exact.] For every fixed $c>0$, the exponential adaptive
schedule $Q_N=e^{cN}$ forces the stated coarse finite-section prime-tail
envelope to zero.

\item[SOH-N006, numerical certificate.] For $w=0.8$, $Q_0=100000$, $c=2$ and
$N\le20$, the sharp v0.5 certificates all pass $10^{-12}$ and have maximum
$3.280365246530569\times10^{-14}$ at $N=5$.
\end{description}

\section{Remaining gap}

The adaptive schedule removes one obstruction: increasing basis size need not
invalidate controlled cutoff removal, provided the cutoff grows with the
basis.  The following remain open:

\begin{enumerate}
\item a useful bound uniform in $N$ at one fixed cutoff;
\item positivity of every infinite-cutoff finite Hermite section;
\item closability and continuity of the complete arithmetic form;
\item identification of its null space with native PhaseNav closure;
\item the full bridge \texttt{SOH-C005}.
\end{enumerate}

Neither the analytic schedule nor its numerical audit proves the Riemann
Hypothesis.
